
\documentclass[12pt,a4paper]{article}

\title{\Huge\bfseries Introduction to LaTeX Document Structure}
\author{
    Siddarth Kumar \\
    \small MNIT Jaipur
}
\date{\today}

\begin{document}

\maketitle

\begin{abstract}
This document demonstrates the basic structure of a
LaTeX document. It covers document classes, title
formatting, and the abstract environment.
\end{abstract}

\section{Introduction}

LaTeX provides different document classes depending
on the type of document being created.

\section{Document Classes}

The document class determines the overall structure
and behavior of a document.

Some commonly used document classes are:

\begin{itemize}
    \item \textbf{article}: Used for short documents,
    papers and reports without chapters.

    \item \textbf{report}: Used for longer documents
    that require chapters.

    \item \textbf{book}: Used for books and theses,
    supporting chapters and front and back matter.
\end{itemize}

\subsection{Choosing a Document Class}

The document class is specified using
\texttt{\textbackslash documentclass}.

For example:

\begin{verbatim}
\documentclass{article}
\documentclass{report}
\documentclass{book}
\end{verbatim}

We can also specify font size and paper size:

\begin{verbatim}
\documentclass[12pt,a4paper]{article}
\end{verbatim}

\section{Customizing the Title}

We can customize the title and author using
formatting commands.

\begin{verbatim}
\title{\Huge\bfseries My Document}
\author{\Large Sahil Kumar}
\date{\today}
\end{verbatim}

Additional information can be placed below the title.

\begin{center}
    \textbf{Supervisor:} Faculty Name \\
    \textbf{Course:} Technical Writing
\end{center}

\section{Abstract}

The abstract environment provides a short summary
of the document.

It is generally placed after the title and before
the main content.

\section{Bibliography and Citations}

BibTeX helps us manage references using a
separate bibliography file.

We can cite a book using the cite command
\cite{kn:gnus}.

\bibliographystyle{plain}
\bibliography{test}
\end{document}
